% ---- pg-540 （书页 528） ----

\[
X_3=\begin{cases}
1, & \text{棕色眼睛}\\
0, & \text{其他}
\end{cases}
\qquad
X_4=\begin{cases}
1, & \text{金色头发}\\
0, & \text{其他}
\end{cases}
\]

\[
X_5=\begin{cases}
1, & \text{右手优势}\\
0, & \text{左手优势}
\end{cases}
\qquad
X_6=\begin{cases}
1, & \text{女性}\\
0, & \text{男性}
\end{cases}
\]

个人 1 和个人 2 在 $p=6$ 个二值变量上的得分为

\begin{mvtable}
\begin{tabular}{c|cccccc}
\toprule
个人 & $\mathbf{X}_1$ & $\mathbf{X}_2$ & $\mathbf{X}_3$ & $\mathbf{X}_4$ & $\mathbf{X}_5$ & $\mathbf{X}_6$\\
\midrule
1 & 0 & 0 & 0 & 1 & 1 & 1\\
2 & 1 & 1 & 1 & 0 & 1 & 0\\
\bottomrule
\end{tabular}
\end{mvtable}

配对与错配的数目由下面的双向表给出.

\begin{mvtable}
\begin{tabular}{ll|cc|c}
\toprule
 & & \multicolumn{2}{c|}{个人 2} & \\
\cmidrule(lr){3-4}
 & & 1 & 0 & 合\quad 计\\
\midrule
个人 1 & 1 & 1 & 2 & 3\\
 & 0 & 3 & 0 & 3\\
\midrule
合\quad 计 & & 4 & 2 & 6\\
\bottomrule
\end{tabular}
\end{mvtable}

用给出等权配对的相似系数 1, 得到
\[
\frac{a+d}{p}=\frac{1+0}{6}=\frac{1}{6}
\]
继续用相似系数 1, 便可得到其余各对个人之间的相似性数字, 这些数字列于下面的 $5\times 5$ 对称矩阵中.

\[
\begin{array}{r@{\;}c}
 & \begin{array}{ccccc}1&2&3&4&5\end{array}\\[2pt]
\begin{array}{r}1\\2\\3\\4\\5\end{array} &
\left[\begin{array}{ccccc}
1&&&&\\
\frac{1}{6}&1&&&\\
\frac{4}{6}&\frac{3}{6}&1&&\\
\frac{4}{6}&\frac{3}{6}&\frac{2}{6}&1&\\
0&\tikz{\node[draw,ellipse,inner sep=1pt]{$\frac{5}{6}$};}&\frac{2}{6}&\frac{2}{6}&1
\end{array}\right]
\end{array}
\]

根据相似系数的大小, 可以断言第二人与第五人最相似, 而第一人与第五人最不相似. 其余各对处于这两个极端之间. 如果我们根据相似性数字将这五个人分成相对接近的两组, 则可得出(1 3 4)和(2 5).

注意, $X_3=0$ 意味着眼球颜色不是棕色的, 故因此一个有蓝眼睛的人与一个有绿眼睛的人组成一个数 0—0 配对. 因此, 用相似性系数 1, 2 或 3 可能不合适, 因为这些系数对 1—1 配对和 0—0 配对赋予相等的权.\hfill$\blacksquare$

我们描述了距离和相似性的构造. 由距离来构造相似性总是可能的. 例如, 我们可以令
\phantomsection\label{eq:12-8}
\begin{equation}
\tilde{s}_{ik}=\frac{1}{1+d_{ik}}
\tag{12-8}
\end{equation}
